% Bagian: second-order-logic
% Bab: metatheory
% Subbagian: undecidability-and-axiomatizability

\documentclass[../../../include/open-logic-section]{subfiles}

\begin{document}

\olfileid[id]{sol}{met}{nax}

\olsection{Logika Orde Kedua Tidak Dapat Diaksiomatisasi}


\begin{thm}
\ollabel{thm:sol-undecidable}
Logika orde kedua tidak dapat diputuskan.
\end{thm}

\begin{proof}
Suatu !!{sentence} orde pertama valid dalam logika orde pertama jika dan hanya
jika kalimat itu valid dalam logika orde kedua, sedangkan logika orde pertama
tidak dapat diputuskan.
\end{proof}

\begin{thm}
\ollabel{cor:sol-not-axiomatizable}
Tidak ada sistem !!{derivation} yang sahih dan lengkap bagi logika orde kedua.
\end{thm}

\begin{proof}
Misalkan $!A$ adalah !!a{sentence} dalam bahasa aritmetika. $\Sat{N}{!A}$
jika dan hanya jika $\Th{PA^2} \Entails !A$. Misalkan $!P$ adalah konjungsi
kesembilan aksioma dari~$\Th{PA^2}$. $\Th{PA^2} \Entails !A$ jika dan hanya
jika $\Entails !P \lif !A$, yakni $\Sat{M}{!P \lif !A}$. Sekarang tinjau
!!{sentence}
$\lforall[z][\lforall[u][\lforall[u'][\lforall[u''][\lforall[L][(!P' \lif !A')]]]]]$
yang diperoleh dengan mengganti $\Obj{0}$ dengan $z$, $\prime$ dengan
variabel fungsi satu-tempat $u$, $+$ dan $\times$ masing-masing dengan
variabel fungsi dua-tempat $u'$ dan $u''$, serta $<$ dengan variabel relasi
dua-tempat~$L$, lalu menguantifikasi semuanya secara universal. Kalimat ini
valid dalam logika orde kedua murni jika dan hanya jika kalimat asal valid,
jika dan hanya jika $\Th{PA^2} \Entails !A$, jika dan hanya jika
$\Sat{N}{!A}$. Jadi, jika terdapat sistem pembuktian yang sahih dan lengkap
bagi logika orde kedua, kita dapat menggunakannya untuk mendefinisikan suatu
enumerasi yang dapat dihitung $f\colon \Nat \to \Sent[L_A]$ atas
!!{sentence}s yang benar dalam~$\Struct{N}$. Fungsi ini dapat
direpresentasikan dalam~$\Th{Q}$ oleh suatu formula orde pertama~$!B_f(x,
y)$. Maka !!{formula}~$\lexists[x][!B_f(x, y)]$ akan mendefinisikan himpunan
!!{sentence} orde pertama yang benar dalam~$\Struct{N}$, bertentangan dengan
Teorema Tarski.
\end{proof}



\end{document}
